\chapter{背景、冲突与上下文}

\section{问题背景}

Pebble 所面对的问题首先来自高压人际互动的普遍性。用户在职场、亲密关系和日常社交中，常常会接触到带有指责、控制、道德绑架或情绪投射色彩的话语。这类话语在表面上可能只是催促、抱怨或询问，但在具体关系中往往还携带更复杂的压力结构。用户若只能凭即时情绪反应作答，回应质量就容易受到恐惧、内疚、愤怒或自我怀疑的影响。

在这一背景下，Pebble 试图处理的并非抽象的文本理解任务，而是一个更贴近真实沟通场景的问题：当用户收到一句令人不适的话时，系统能否帮助其先分辨真实含义，再形成较为稳健的回应，并在后续互动中持续积累更适配的策略。读心翻译器、模拟陪练场和急救呼吸三个功能原型都围绕这一目标展开，只是它们分别切入了理解、练习与稳定状态的不同阶段。

\section{核心冲突}

本项目的核心冲突在于，用户需要的是带关系上下文的沟通支持，而多数传统交互系统提供的是脱离上下文的单轮回答。若系统只处理当前这句话，它通常能够生成语义层面的解释，却难以判断这句话在特定关系中究竟意味着催促、试探、操控、情绪投射还是边界侵犯。对于 Pebble 这类产品而言，真正有价值的判断并不来自单句文本本身，而来自文本、关系历史、场景压力与用户目标之间的联动。

这一冲突进一步表现为两个层面的张力。

\begin{itemize}
  \item 用户希望快速获得帮助，因此系统需要在高压时刻给出低摩擦、可执行的支持。
  \item 用户又希望系统足够懂自己所处的关系背景，因此系统不能长期停留在一次性分析层面。
\end{itemize}

前者要求产品具备即时性，后者要求系统具备持续记忆与情境适配能力。Pebble 的架构与产品设计正是在这两类要求之间寻找可行平衡。

\section{上下文约束}

Pebble 所处的上下文并不只是“做一个 AI 页面”。从产品侧看，系统必须同时处理至少三类不同但彼此相关的需求。

\begin{itemize}
  \item 即时解释：用户在收到令人不适的话语后，希望快速获得潜台词分析与候选回复。
  \item 安全练习：用户希望在真实冲突发生前进行模拟对话，观察自己是否仍然陷入辩解、讨好或情绪升级。
  \item 状态恢复：在用户已被情绪淹没时，系统需要提供比分析更基础的稳态支持。
\end{itemize}

从系统侧看，Pebble 还面临两类结构性约束。

\begin{itemize}
  \item 关系上下文必须可持续维护，否则同一句话在不同关系中的差异无法被稳定表达。
  \item 系统边界必须保持克制，否则产品很容易滑向心理诊断、人格贴标签或过度承诺的方向。
\end{itemize}

也正因如此，Pebble 在相关设计文档中强调去标签化、渐进累积、用户主权与隐私优先。这些约束决定了本报告后续不能把产品简单写成“更聪明的情绪识别器”，而必须把它写成一个围绕关系情境组织支持链路的系统。

\section{本报告的回答方式}

针对上述背景与冲突，本报告给出的回答方式是把 Pebble 解释为一个关系上下文驱动的闭环系统。第一章和第二章已经分别确立了项目定义与主张；从本章开始，报告将进一步说明这一主张为何成立。具体来说，后续章节会依次回答四个问题：为什么单句分析不足以支撑高质量沟通支持，Pebble 需要哪些核心概念与设计原则，系统如何把解码、陪练、稳态恢复与长期记忆组织成可协同的结构，以及当前产品骨架与实现边界处在什么位置。

因此，本章的作用是把读者带到一个更准确的问题入口，而不提前展开最终设计细节。后文结构建立在以下判断之上：Pebble 所讨论的是如何在复杂关系中为用户建立一个可持续的理解、练习与应对支撑系统。
